\chapter{Discussion}
\label{ch:discussion}

\section{What the degradation architecture bought, and what it did not}

The premise in \cref{ch:intro} was that a body-worn master will lose channels,
so losing them must be survivable. The two dated captures support the premise
empirically: the failed population grew between them
(\cref{sec:channelhealth}). The architecture's value can be stated precisely
from the measurements.

It kept the system running. With eight of fourteen channels failing the health
criterion, the right arm retains full position sensing at $R^2 = \num{0.986}$
and the left arm retains direction. A joint-space mapping would have been
inoperable on either arm from the first failed channel, because it has no
reconstruction path.

It kept the failures visible. The incoherent verdict exists because a failing
potentiometer does not go quiet, and range alone calls a wide-spanning
nonsense signal healthy. Without that verdict the system would have kept
following \SI{54.5}{\milli\metre} of commanded motion the operator never
asked for, on a master whose whole reach is \SI{272}{\milli\metre}.

What it did not buy is capability. Freezing the left arm's bend joints removes
its radial dimension entirely, collapsing its commanded set from a shell to a
surface, and the wrist-roll fallback restores an axis rather than restoring
the measurement. The architecture converts a fault into a known, bounded,
announced reduction. It does not convert it into working hardware, and
\cref{sec:channelhealth} names the two channels whose repair would.

\section{The reduced-degree-of-freedom premise, assessed}

The second premise was that commanding fourteen joints through a body-worn
master is a usability problem, and that the answer is to take each component
from the sensor that observes it best. The evidence is mixed, and it divides
cleanly by component.

Elevation from gravity works. It is mount-independent, drift-free, integrates
nothing, and survived every reformulation. Reach from the forward-kinematic
magnitude works, because a magnitude depends on the bend joints and tolerates
a dead roll passed through as zero. Both vertical and fore-and-aft motion
track correctly on both arms.

Azimuth from a single joint does not work, and the reason is instructive
rather than incidental. The spherical model assumes azimuth and elevation are
independent. They are not, because joint one is a roll about the arm's own
axis, so its effect on the tip scales with how bent the arm is. The
\SI{137.7}{\degree} measured between two directions intended to be
\SI{90}{\degree} apart is not a calibration error that a better axis map or a
better zero could remove; it is the model being wrong about the geometry.

The honest reading is that the decomposition succeeded for the two components
where a single sensor observes the quantity directly, and failed for the one
where it does not. That is a usable design rule for a successor: decompose
along the axes your sensors actually observe, not along the axes that make the
mathematics tidy.

\section{The workspace result and what it costs the study}

\Cref{ch:workspace} reports that the two arms share no reachable workspace
under the as-built mount. The consequence for the experimental programme is
severe and worth stating plainly rather than absorbing into a limitations
list.

Inter-arm handover is impossible. Not difficult, not unreliable: there is no
point in space both arms can reach, so no control strategy, no planner and no
orientation mode can produce one. The task that would have demonstrated
autonomy achieving something teleoperation cannot was blocked by geometry
before it was blocked by anything else, and the finding is the reason it is
not in the task set.

The cause is not the mount rotation, which was proved to leave the mirror
residual invariant. It is that the two arms are parked asymmetrically at their
home joint angles, and those angles are ground truth because the simulation to
hardware bridge replays them onto the real arm. Changing them in software
would command the difference as a jump. Fixing this requires re-parking the
right arm in hardware.

Two of the five delivered tasks are the same bimanual mechanism with different
objects, and only one task uses the second mechanism. There are two genuinely
bimanual mechanisms available on this platform, not three. That is a real
narrowing of what the platform can be used to ask.

\section{Comparison against the planning report's targets}

The delivered system diverges substantially from the plan, and the divergences
fall into three groups.

\emph{Superseded by measurement.} The planning report specified physical
springs at the master for force feedback and virtual spring impedance at the
slave. The impedance target is unreachable on this platform for a reason
discovered during the work: the arms' cyclic control interface is unusable
over the available network path, so the system commands through a high-level
velocity interface at \SIrange{18.4}{18.7}{\hertz}. An impedance law at
\SI{18}{\hertz} is not an impedance law. The compliance the plan sought would
have to come from the mechanism instead, through series elastic actuation at
the joint~\cite{toubar2022sea,lee2019sea}, or from a current-based
formulation that does not require a high-rate torque
loop~\cite{dewolde2024impedance}. Both are hardware changes.

\emph{Deferred and still open.} Coordinated pick-and-place on the real
platform was not demonstrated. The task set is specified and verified
geometrically, and the runner exists, but the arms were not available for the
period in which it would have been run.

\emph{Added.} The safety architecture, the degradation architecture and the
workspace characterisation were not in the plan. The first two are the direct
consequence of discovering that sensor faults dominate; the third is the
consequence of asking whether the task set was reachable before running it.

\section{The safety layer is thorough about distance and silent about time}

The safety architecture is careful about the quantity it was designed around.
The margin is measured geometrically against the parts of the arm that
actually pass the wearer rather than through the semantic collision
description; fourteen injected faults are caught; every mechanism that can
stop the arm announces itself by name; and a blocker that cannot state how it
clears is refused at registration. Most of that was built after a specific
failure demonstrated the need for it.

None of it addresses the fact that the margin is a \emph{distance} and the
wearer can move. \Cref{tab:reaction} gives the arithmetic: in five of six
modelled limb motions the limb travels further than the entire
\SI{150}{\milli\metre} margin before the guard is aware anything has changed.
The system protects a stationary person from a moving robot, and it is
deployed on an arrangement that also contains a moving person.

Three responses are available and none is implemented. Inflating the floor by
the perception latency multiplied by an assumed limb speed would at least make
the margin mean what it claims, and it is a one-parameter change. Cutting the
latency is largely available, since most of the \SI{562}{\milli\second}
spread is a configured publication rate limit rather than the estimator, which
runs at \SIrange{28}{34}{\milli\second} per frame. And scaling commanded speed
with measured clearance would couple two quantities the system already
computes and currently keeps apart. The first is the one to do first, because
it changes what the margin asserts rather than merely improving it.

\subsection{The fallback restores an assumption that is optimistic by
\SI{224}{\milli\metre}}

Admitting a camera into the clearance check is governed by the monotonicity
rule of \cref{sec:wearer-tracking}, and that rule is correct about its own
scope: a tracking result that would place the wearer further away changes
nothing. The gap is on the other side. When the estimate is stale, silent or
gated out, the guard falls back to the mannequin, whose arms are down.

Measured through the guard's own geometry, a wearer holding their arms out to
the sides has \SI{0.0933}{\metre} of clearance at the home pose, where the
fallback believes there is \SI{0.3170}{\metre}. The system's present answer to
this is an instruction to the wearer, recorded as a standing operating rule.
That is a genuine mitigation and it is not a control system, and it should be
read alongside the reaction budget above rather than separately: markerless
tracking on this platform is load bearing, and its failure mode is to restore
an assumption quietly.

\subsection{And the clearance model itself is not yet trustworthy}

The most uncomfortable single finding of the hardware session belongs here
rather than in the results. Asked about a pose in which the arm was
\emph{physically against the mannequin}, the clearance model returned
\SI{367.0}{\milli\metre}, which is the same value it returns for a visibly
clear pose. It samples link \emph{origins} rather than link geometry, so it
cannot see the arm between joints.

Every clearance figure in this thesis that was computed through the
guard's capsule model rather than through that link-origin model stands, and
they are the majority. But the check that runs on the physical arm during
homing was the link-origin one, which means the standing constraint that the
wearer's margin is enforced at all times is, at the time of writing,
\textbf{not enforced in code on hardware}. \Cref{ch:conclusion} places its
repair third in the order of work, above the mount change and above the study.

\section{What the perception stage did and did not change}

Replacing declared coordinates with measured ones removes a specific and
serious failure mode: an arm acting with complete confidence on a number in a
file that no longer describes the room. The map is accurate enough for the
task, at every object centre within \SI{1}{\milli\metre}, and it is the first
structure in this project a planner has consulted, which closes the gap in
which a path could route correctly around the wearer and sweep through the
bench.

It does not make the platform able to pick things up. The hardware session in
\cref{sec:hardware} found that the transform from the camera frame to the arm
is wrong by roughly \SI{8.5}{\degree} and rotates with the wrist, which is
larger than the entire capture gate at working distance. A perfect map
expressed through a wrong extrinsic yields a wrong grasp. The correct reading
is that the perception work moved the bottleneck from \emph{not knowing where
things are} to \emph{not knowing where the camera is}, and the second is a
calibration with a known method, which is a better position to be in.

\subsection{What the real cameras changed about that assessment}

\Cref{ch:vision} is the first evidence that the geometry behaves on real
sensors, and it changes the assessment in three specific ways.

\emph{The surface-fitting stage transfers.} A support plane at
\SI{1.43}{\milli\metre} root mean square with \SI{74}{\percent} inlier support
at a working distance of \SI{0.278}{\metre}, on a real table under laboratory
lighting, is within a factor of three of what the same code achieves against a
renderer. That is the one component of the perception chain that can now be
said to work outside simulation, and it is the component the pick path leans
on hardest.

\emph{The scene camera is a different instrument and should not be used as a
work-surface sensor.} At a median range of \SI{4.7}{\metre} it returns depth
on \SI{14.5}{\percent} of its pixels and its best-supported plane holds
\SI{3.8}{\percent} of the points it does return. The design documents treated
it as a wide-angle equivalent of a wrist camera; it is not, and every quantity
that depends on a support plane must come from a wrist camera. This is a
limitation of the arrangement rather than of the device: there is no dominant
plane in a room.

\emph{The learned components contributed nothing that the classical ones did
not.} The open-vocabulary detector returned no instance of the prompt it was
given on any of the four cameras, while the colour-and-geometry path found the
objects and measured them. This is a single scene and a single prompt, so it
is not a general claim about the models; but it does mean that no result in
this thesis rests on a learned detector, which is worth knowing given that
\cref{ch:method} argues the models were adopted for open vocabulary rather
than for accuracy.

The uncomfortable part of \cref{ch:vision} is not any of those. It is that
pointing an instrumented pipeline at the hardware for the first time
immediately found a constant that three files asserted, none measured, and
which had silently inverted every scene figure the project had recorded. The
defect was undetectable by any metric the pipeline computed, because a
consistently applied rotation preserves every distance. It was detectable only
by comparing one camera's picture against another's. The general lesson is the
one \cref{ch:instrument} has been accumulating: the quantities most likely to
be wrong are the ones nobody thought were quantities.

\section{Validity of the measurements}
\label{sec:validity}

\subsection{The base rate of instrument error, and what caught it}

Twenty-six times during this project a measurement was wrong and the system
was not. That is a high enough base rate to change how the rest of this
chapter should be read, so it is stated as a result rather than left in a
methods note.

The pattern across the twenty-six is more useful than any individual case.
Sorted by mechanism they fall into five groups: an absence read as a value; a
test harness that constructs the condition it is testing; state left over from
a previous run; the order of operations inside the measurement itself; and ---
added by this session --- a constant asserted in more than one place and
measured in none. None is specific to robotics.

Two numbers carry the argument. \textbf{Fifteen of the twenty-six were caught
by an independently known quantity contradicting the reading}; five by
inspecting the code, and six by looking at the raw artefact, usually an image.
Care was not the defence that worked. Redundancy was. And \textbf{fourteen of
the twenty-six were already answerable from something in hand} --- a
published stroke, an arithmetic prediction, a quantity that is exact by
construction --- which is the recommendation this project's experience
supports: build the cross-check in rather than trying harder.

The consequence for this thesis is specific and is applied throughout. Every
reachability figure is reported with its repeat count because one solver call
is one coin flip. Every stage of the simulation campaign of
\cref{sec:campaign} carries controls that must \emph{fail}, and
\cref{sec:sim-cv} reports those alongside the results they license. Where a
quantity has never been measured, the text says so rather than offering an
estimate --- the camera mounting transform in \cref{fig:budget} is drawn with
no bar for exactly that reason. And any unreproduced single measurement in
this document should be read as provisional.

\Cref{app:instrument} is the full record: all twenty-six cases, what each
actually was, and what revealed it.

\subsection{Simulation is not the arm}

Every quantitative result in this thesis except the channel measurements comes
from simulation with mock hardware. The mock echoes commanded positions with
no dynamics, which has one consequence that must not be glossed: every force
reading is identically zero, so any quantity derived from a force is
structurally undefined rather than approximated. The measurement that payload
configuration made exactly \SI{0}{\metre} of difference to tracking is a
statement about the mock.

\subsection{Latency figures are lower bounds}

Every latency reported is control-path latency, measured between software
stages. It excludes the operator's reaction, the display, and the physical
arm's response. Any end-to-end figure will be larger, and the reported values
should be read as a floor.

\subsection{The one-stack constraint}

Two instances of the master node split the serial stream between them and
invalidated a full day of measurements before the cause was found. This is now
prevented at the file-descriptor level, but it means that any measurement
taken before that prevention was added carries a risk that cannot be
retrospectively excluded.

\section{The study the platform is not ready for}
\label{sec:not-run}

The two-person study this architecture was built to enable was designed in
full and was not run. Both halves of that sentence are results, and they are
reported here rather than as a chapter of their own, because a protocol that
has not been executed is not a finding about supernumerary limbs.

What was produced is a complete protocol (\cref{app:study}) and a measured
property of the analysis that would evaluate it (\cref{app:analysis}). The
second is worth stating in the body. Applied to \num{20000} constructed
studies at the planned sixteen dyads, the pre-registered procedure recovers
the effect the design was powered for on \SI{97.3}{\percent} of them
family-wise, but on only \SI{63.8}{\percent} for any \emph{named} task once
the correction is applied, while rejecting a true null at exactly the nominal
rate. The design's own power calculation reports neither figure. So the study
as designed is well powered for the question it asks --- whether shared
autonomy changes performance on this platform --- and under-powered for a
per-task claim, and a per-task null at $n = 16$ would have to be reported as
an interval rather than as an absence.

It was not run for reasons that are specific rather than administrative.

\begin{enumerate}
  \item \textbf{Two master channels.} Left j2 and left j4. With both frozen
    the left arm has no radial dimension, so the direct teleoperation
    condition, which is the baseline the whole comparison rests on, is not
    representative of a working master.
  \item \textbf{The right arm's home pose.} It has never been read from
    hardware, and the workspace geometry depends on it.
  \item \textbf{Hardware validation of the safety layer.} Nothing in the
    safety architecture has run against a real arm. The Kortex session
    recovery path is the riskiest part, because the arm permits exactly one
    session and a leaked one refuses the next connection.
  \item \textbf{Ethics approval for a two-person protocol}, which differs from
    a single-participant protocol in that one participant bears a physical
    risk arising from another participant's actions.
  \item \textbf{A detection rate measured on real images.} The synthetic
    figure characterises the renderer.
  \item \textbf{A hand-eye calibration for the wrist camera.} The transform is
    wrong by roughly \SI{8.5}{\degree} and is the largest term in the
    positioning budget. It needs the hardware and a calibration target, and
    until it exists no open-loop grasp will land inside the capture gate.
  \item \textbf{A clearance check that measures link geometry rather than link
    origins}, before anybody wears the rig. This is the one blocker on the
    list that is a safety defect rather than an incompleteness.
\end{enumerate}

None of these is a research question. All are execution, and the first is two
soldered joints.
